%! @@TITLE@@ — Unit @@UNITINT@@, Lesson @@LESSONID@@ — Guided Notes
% GRADUAL RELEASE: this is the "I do / we do" block — 15 minutes. The vocabulary is BUILT HERE,
% before the group activity uses it. Reference implementation: unit01/lesson00/notes.
\documentclass[10pt]{article}
\usepackage{@@PREFIX@@-article}
\usepackage{@@PREFIX@@-boxes}
\newcommand{\TallMath}[1]{$\displaystyle #1\rule[-1.4em]{0pt}{3.2em}$}
% Coordinate grid: {xmin}{xmax}{ymin}{ymax}
\newcommand{\lgrid}[4]{%
  \draw[step=1, linegray!40, very thin] (#1,#3) grid (#2,#4);
  \draw[->, semithick, charcoal] (#1,0)--(#2,0) node[below right, font=\scriptsize]{$x$};
  \draw[->, semithick, charcoal] (0,#3)--(0,#4) node[left, font=\scriptsize]{$y$};}

\begin{document}

\pageheader{Unit @@UNITINT@@, Lesson @@LESSONID@@}{Guided Notes}
% NAMESTRIP: no \namedateperiod here. The student writes their name once, on the cover.

\vspace{0.10in}

\begin{objectivebox}
\textbf{By the end of this lesson, I will be able to\dots}
\begin{itemize}\setlength{\itemsep}{2pt}
  \item TODO: objective 1 --- use the formal vocabulary freely; there is no spoiler rule.
  \item TODO: objective 2.
\end{itemize}
\end{objectivebox}

\vspace{0.05in}

\begin{vocabbox}
\small
Fill in each term as we build it together.
\par\vspace{2pt}   % \par is required: \termblanklong's \noindent is a no-op mid-paragraph
\termblanklong{TODO term 1}
\termblanklong{TODO term 2}
% 5–7 terms. The key defines a matching \vocabans macro — see notes_key.tex.
\end{vocabbox}

\begin{hookbox}
TODO: a 60-second context that motivates the lesson and supplies the numbers every worked example
below will reuse. Ask 3--4 quick questions with fill-in answers, then land the idea in one
sentence \emph{before} naming anything.
\begin{itemize}\setlength{\itemsep}{1pt}
  \item TODO: \blank{1.0cm} \quad \blank{1.0cm} \quad \blank{1.0cm}
\end{itemize}
\end{hookbox}

% FOUR numbered sections. Keep ONE worked context running through all of them where you can —
% it is what makes the block read as one lesson rather than four procedures.
\boxguard[24]
\begin{notesbox}{1. TODO: the defining idea}
TODO. Build the definition from the hook's own numbers, then state the general form.
\begin{work}
  % TODO: the first worked example, one statement per line, & before the relation.
  m &= \frac{y_2 - y_1}{x_2 - x_1} \\
    &= 2
\end{work}
\end{notesbox}

\boxguard[22]
\begin{notesbox}{2. TODO: the second idea}
TODO --- a fill-in \texttt{tabularx} works well here (feature / how to find it / what it means).
\end{notesbox}

\boxguard[26]
\begin{notesbox}{3. TODO: the section that carries the target misconception}
TODO. Put the two things students conflate side by side in a two-column \texttt{tabularx}, then
close with a gold caution box giving a case where the two answers \emph{disagree}:
\begin{tcolorbox}[colback=goldbg, colframe=goldacc, boxrule=0.8pt, arc=1.5mm,
    left=2.5mm, right=2.5mm, top=1.5mm, bottom=1.5mm]
\textbf{\textcolor{goldacc}{Careful: TODO.}} \; TODO --- the counterexample, with fill-in blanks.
\end{tcolorbox}
\end{notesbox}

\boxguard[20]
\begin{notesbox}{4. TODO: the edge cases and the boundary}
TODO --- the special case, the value where the rule fails, and the story-vs-math distinction.
\end{notesbox}

% GUIDED PRACTICE — the "we do". Everything at once on ONE new example; do item 1 together, then
% release the rest and circulate.
\begin{practicebox}
TODO: one new example, all features at once.
\begin{enumerate}[leftmargin=1.7em, itemsep=3pt]
  \item TODO \blank{1.2cm}
  \item TODO --- solve it here.
        \begin{work}
          % TODO
          2x - 4 &= 0 \\
               x &= 2
        \end{work}
\end{enumerate}
\end{practicebox}

\end{document}
